\section{Conclusion and Limitations}
\label{sec:conclusion}

We presented MAW, a framework that connects orchestration-task construction with feedback on how delegated work is executed. Recursive composition expands functional goals within existing environments; semantic reconstruction and execution-grounded validation turn these candidates into training tasks. An outcome-prioritized reward combines answer and database-state verification with tool progress, answer coverage, and collaboration feedback to train an orchestrator with frozen workers. The evaluation separates complete-system transfer, orchestration benefit, and training-strategy effects. \drafttodo{Add the principal measured findings after the six-benchmark evaluation and controlled ablations are complete.}

The method has several limits. Reference matching may under-credit correct alternative solutions, and independent-goal rewards apply only to a verified subset of tasks. Logical model rounds omit tool latency and generation length, so efficiency claims require separate wall-clock measurements. Offline recursive construction does not establish online co-evolution of environments and policies. The audited two-parent synthesis snapshot also does not establish the benefit of multiple recursive rounds. The current interface requires delegation, and evidence from one model scale with frozen workers cannot establish generality to unrestricted agent organizations or jointly trained teams. Generated environments and completion checks also require audits for semantic errors that successful replay may not expose.
